% Para utilizar este template, consulte o README.md e os guias de normalização
% publicados pelo Sistema de Bibliotecas da Universidade Federal do Ceará.
%
% Atualização normativa e técnica: Tiago Guimarães Sombra (2026).
% Versão comunitária: 1.1.1 (2026-08-17).
% Base preservada do template UFC/abnTeX2 2016/2019.
%
% Este arquivo concentra a infraestrutura LaTeX comum. Regras de apresentação
% específicas da UFC e a interface \ufcsetup ficam em lib/ufctex.sty.
%
% PRINCÍPIO DA MODERNIZAÇÃO:
% - manter abnTeX2/memoir como base estável e retrocompatível;
% - usar APIs atuais do LaTeX quando possível;
% - substituir dependências legadas somente quando houver ganho verificável;
% - manter recursos pesados (gráficos, minted, caixas, pseudocódigo avançado)
%   como módulos opcionais carregados por \ufcsetup antes de \begin{document}.

% -----------------------------------------------------------------------------
% Engine, codificação e tipografia
% -----------------------------------------------------------------------------
% LaTeX moderno assume UTF-8 por padrão; inputenc deixa de ser necessário.
% iftex permite que o mesmo projeto compile com pdfLaTeX, LuaLaTeX ou XeLaTeX.
\usepackage{iftex}

\ifPDFTeX
  \usepackage[T1]{fontenc}
  % Times-compatible estável e de alta qualidade no fluxo pdfLaTeX.
  \usepackage{newtxtext,newtxmath}
\else
  % Fluxo Unicode/OpenType para LuaLaTeX/XeLaTeX.
  \usepackage{fontspec}
  \defaultfontfeatures{Ligatures=TeX}
  % O TeX Live/Overleaf distribui a variante OpenType TermesX com newtx.
  % Usar os nomes dos arquivos evita dependência do cache de fontes do sistema.
  % Caso uma unidade exija literalmente Times New Roman e a fonte esteja
  % legalmente instalada no ambiente, substitua esta configuração.
  \setmainfont{TeXGyreTermesX-Regular.otf}[
    BoldFont=TeXGyreTermesX-Bold.otf,
    ItalicFont=TeXGyreTermesX-Italic.otf,
    BoldItalicFont=TeXGyreTermesX-BoldItalic.otf
  ]
\fi

% LEGADO:
% \usepackage[utf8]{inputenc} % desnecessário no LaTeX atual
% \usepackage{mathptmx}       % substituído por newtx no fluxo pdfLaTeX

% COMPATIBILIDADE OVERLEAF / TEX LIVE 2025:
% Algumas combinações do microtype distribuído no TeX Live 2025 e do kernel
% LaTeX 2025 emitem o falso-positivo "Command \\showhyphens has changed".
% Filtramos somente essa mensagem conhecida, preservando os demais warnings.
\usepackage{silence}
\WarningFilter{latex}{Command \showhyphens has changed}
\usepackage{microtype}

% -----------------------------------------------------------------------------
% Matemática e unidades
% -----------------------------------------------------------------------------
\usepackage{mathtools} % inclui e estende amsmath
\ifPDFTeX
  % newtxmath já foi carregado junto da tipografia no fluxo pdfLaTeX.
\else
  % unicode-math deve ser carregado depois de mathtools. As duas mensagens
  % abaixo documentam redefinições intencionais de comandos compartilhados
  % entre os pacotes e são filtradas pela própria API de unicode-math.
  \usepackage[warnings-off={mathtools-overbracket,mathtools-colon}]{unicode-math}
  \setmathfont{XITSMath-Regular.otf}
\fi
\usepackage{siunitx}
\sisetup{
  output-decimal-marker = {,},
  group-separator = {.},
  group-minimum-digits = 4,
  mode = match,
  propagate-math-font = true,
  reset-math-version = false,
  reset-text-family = false,
  reset-text-series = false,
  reset-text-shape = false,
  text-family-to-math = true,
  text-series-to-math = true
}

% -----------------------------------------------------------------------------
% Figuras e objetos flutuantes
% -----------------------------------------------------------------------------
% Política única de separação visual entre identificação, objeto e fonte.
\newlength{\UFCCaptionObjectSep}
\setlength{\UFCCaptionObjectSep}{6pt}
\newlength{\UFCObjectSourceSep}
\setlength{\UFCObjectSourceSep}{6pt}

\usepackage{graphicx}
\graphicspath{{figuras/}}
\usepackage[export]{adjustbox}
\usepackage{flafter}
\usepackage{placeins}
\usepackage[
  font={footnotesize,singlespacing},
  format=plain,
  justification=justified,
  singlelinecheck=false,
  skip=6pt
]{subcaption}
\usepackage{pdflscape}

% -----------------------------------------------------------------------------
% Tabelas e quadros
% -----------------------------------------------------------------------------
% Compatibilidade com documentos antigos: os ambientes tradicionais continuam
% disponíveis. Para tabelas novas, o template também fornece tabularray-abnt,
% que usa uma implementação LaTeX3 e possui temas específicos para tabela/quadro.
\usepackage{array,tabularx,longtable}
\usepackage{booktabs}
\usepackage{tabularray-abnt}
\UseTblrLibrary{booktabs,siunitx}
\SetAbntTblrFont{\normalsize}
% Captions internas de tall/long tabularray também devem seguir a UFC:
% fonte 10, entrelinha simples e separação visual de 6 pt.
\SetTblrStyle{caption,lasthead,capcont}{font=\footnotesize}
\SetTblrStyle{firsthead-text,lasthead-text,conthead-text}{font=\footnotesize}
\SetTblrStyle{lastfoot}{font=\footnotesize}

% LEGADO DISPONÍVEL SOMENTE SE O USUÁRIO PRECISAR:
% multirow, xltabular, threeparttable e makecell deixaram de ser carregados por
% padrão porque tabularray cobre os casos novos e reduz sobreposição de pacotes.
% Documentos antigos que dependam diretamente deles podem reativá-los aqui:
% \usepackage{multirow}
% \usepackage{xltabular}
% \usepackage{threeparttable}
% \usepackage{makecell}

% -----------------------------------------------------------------------------
% Texto, código-fonte e cores
% -----------------------------------------------------------------------------
\usepackage{indentfirst}
\usepackage{xcolor}
\usepackage{listings} % padrão portátil, sem dependências externas
% fvextra é carregado antes de csquotes para manter o módulo opcional minted
% livre de warnings de ordem de carregamento (lineno/csquotes).
\usepackage{fvextra}

% minted e pseudocódigo moderno são módulos opcionais de \ufcsetup.
% LEGADO: algorithm2e também pode ser ativado por \ufcsetup para documentos
% antigos, mas não é mais carregado globalmente.

% -----------------------------------------------------------------------------
% Glossários, índice e elementos auxiliares
% -----------------------------------------------------------------------------
% glossaries-extra estende glossaries e mantém compatibilidade com os comandos
% usuais (\newacronym, \gls, \printglossary), oferecendo infraestrutura atual.
% tracklang pode produzir mensagens de recurso de locale que não afetam a saída;
% a supressão abaixo é restrita ao carregamento do glossário.
\usepackage{tracklang}
\TrackLangShowWarningsfalse
\usepackage[nogroupskip,nonumberlist,translate=babel]{glossaries-extra}
\TrackLangShowWarningstrue

\usepackage{pdfpages}
\usepackage{xurl}
\usepackage{csquotes}

% -----------------------------------------------------------------------------
% Legendas e espaçamento dos objetos acadêmicos
% -----------------------------------------------------------------------------
% Os guias UFC exigem identificação acima e fonte abaixo para ilustrações
% (incluindo figuras, gráficos, fluxogramas e quadros) e tabelas, com fonte 10 e
% espaço simples. Não há distância numérica prescrita entre identificação,
% objeto e fonte. O template adota 6 pt como política tipográfica uniforme.
% Algoritmos e códigos-fonte não recebem regra específica nos guias UFC; quando
% legendados, o template estende a eles a mesma política visual por consistência.
\usepackage[
  font={footnotesize,singlespacing},
  format=plain,
  justification=justified,
  labelsep=endash,
  singlelinecheck=false
]{caption}
\captionsetup{skip=\UFCCaptionObjectSep}
% Subfiguras/subtabelas não são tratadas separadamente pelo guia UFC; quando
% usadas, preservamos a mesma fonte 10 e espaço simples para não criar uma
% legenda menor que o padrão institucional.
% longtblr é a infraestrutura usada por longabnttblr. Seus separadores padrão
% já são 6 pt, mas os fixamos explicitamente para manter uma única política.
\SetTblrOuter[longtblr]{headsep=\UFCCaptionObjectSep,footsep=\UFCObjectSourceSep}
\SetTblrOuter[talltblr]{headsep=\UFCCaptionObjectSep,footsep=\UFCObjectSourceSep}

% -----------------------------------------------------------------------------
% Citações e referências
% -----------------------------------------------------------------------------
% O backend atual é Biber com biblatex-abnt. O template mantém os comandos
% \cite e \citeonline do modelo histórico como aliases de compatibilidade.
\usepackage[
  backend=biber,
  style=abnt,
  language=auto,
  autolang=other,
  maxcitenames=3,
  mincitenames=1,
  maxbibnames=99,
  giveninits=false,
  uniquename=false
]{biblatex}
\addbibresource{3-pos-textuais/referencias.bib}

\let\ufcBiblatexCite\cite
\let\cite\parencite
\ProvideDocumentCommand{\citeonline}{o m}{%
  \IfNoValueTF{#1}{\textcite{#2}}{\textcite[#1]{#2}}%
}

% LEGADO:
% \usepackage[alf,...]{abntex2cite}

% -----------------------------------------------------------------------------
% Biblioteca UFC e configuração institucional
% -----------------------------------------------------------------------------
\usepackage{lib/ufctex}

% -----------------------------------------------------------------------------
% Normalização transversal de objetos
% -----------------------------------------------------------------------------
% A camada histórica ufctex possui implementações próprias para fonte, código e
% alguns ambientes de algoritmo. As regras abaixo centralizam a política desta
% versão sem exigir alterações manuais em cada figura, tabela, quadro ou código.
% A composição do objeto já fornece a separação antes de Fonte/Nota; não se soma
% outro \vspace de 6 pt aqui, evitando duplicar esse afastamento visual.
\renewcommand{\Fonte}[2][\fontename]{%
  \par\noindent%
  {\footnotesize\SingleSpacing #1\ABNTEXcaptionfontedelim #2\par}%
}
\renewcommand{\Nota}[2][\notaname]{\Fonte[#1]{#2}}

% algorithm/algpseudocodex usa o mecanismo do package caption.
\AddToHook{package/algorithm/after}{%
  \captionsetup[algorithm]{%
    font={footnotesize,singlespacing},
    labelfont=normalfont,
    textfont=normalfont,
    format=plain,
    justification=justified,
    labelsep=endash,
    singlelinecheck=false,
    skip=\UFCCaptionObjectSep
  }%
}

% algorithm2e possui mecanismo de legenda próprio.
\AddToHook{package/algorithm2e/after}{%
  \SetAlCapFnt{\footnotesize\linespread{1}\selectfont}%
  \SetAlCapNameFnt{\footnotesize\linespread{1}\selectfont}%
  \SetAlCapSty{textnormal}%
  \SetAlCapNameSty{textnormal}%
  \SetAlgoCaptionSeparator{\space\textendash\space}%
  \SetAlCapHSkip{0pt}%
  \SetAlCapSkip{\UFCCaptionObjectSep}%
  \setlength{\interspacealgoruled}{\UFCCaptionObjectSep}%
}

% listings é carregado no núcleo; a configuração global de caption já
% fornece fonte reduzida e espaço simples. Aqui fixamos apenas posição e gaps.
\lstset{%
  captionpos=t,
  abovecaptionskip=0pt,
  belowcaptionskip=\UFCCaptionObjectSep
}

% Referências cruzadas. cleveref foi mantido por compatibilidade e maturidade;
% zref-clever foi avaliado como alternativa, mas não oferece ganho suficiente
% para justificar quebra da API \cref/\Cref nesta versão do template.
\usepackage{bookmark}
\usepackage[capitalise,noabbrev,nameinlink]{cleveref}
\crefname{figure}{Figura}{Figuras}
\crefname{table}{Tabela}{Tabelas}
\crefname{quadro}{Quadro}{Quadros}
\crefname{equation}{Equação}{Equações}
\crefname{chapter}{Capítulo}{Capítulos}
\crefname{section}{Seção}{Seções}
\crefname{algorithm}{Algoritmo}{Algoritmos}

% Referências: margem esquerda, espaço simples e uma linha simples entre entradas.
\setlength{\bibhang}{0pt}
\setlength{\bibitemsep}{\baselineskip}
\AtBeginBibliography{\SingleSpacing\raggedright}

\makeglossaries
\makeindex
